\documentclass[11pt, a4paper]{article}

% --- Paquetes Fundamentales ---
\usepackage[utf8]{inputenc}
\usepackage[T1]{fontenc}
\usepackage[spanish,es-tabla]{babel}
\usepackage[margin=2.5cm, headheight=15pt]{geometry}
\usepackage{graphicx}
\usepackage{fancyhdr}
\usepackage{xcolor}
\usepackage{microtype}
\usepackage[colorlinks=true, linkcolor=ucbblue, urlcolor=ucbblue]{hyperref}

% --- Matemáticas y Electrónica ---
\usepackage{amsmath, amssymb}
\usepackage{siunitx}
\usepackage[american, RPvoltages]{circuitikz}

% --- Elementos de Diseño Pedagógico y Código ---
\usepackage{tcolorbox}
\usepackage{enumitem}
\usepackage{listings}
\usepackage{tabularx}
\usepackage{booktabs}
\usepackage{colortbl}

% --- Configuración de Colores y siunitx ---
\definecolor{ucbblue}{RGB}{0, 51, 102}
\definecolor{codegreen}{rgb}{0,0.6,0}
\definecolor{codepurple}{rgb}{0.58,0,0.82}
\definecolor{ltspice}{RGB}{128, 0, 0}
\sisetup{
    output-decimal-marker = {,},
    per-mode = symbol,
    inter-unit-product = \ensuremath{{}\cdot{}}
}

% --- Macros y Estilos ---
\newcommand{\tecla}[1]{\colorbox{gray!20}{\texttt{\textbf{#1}}}} % Macro para atajos de teclado

\tcbset{
    caja_diagnostico/.style={colback=blue!3, colframe=ucbblue, fonttitle=\bfseries, boxrule=0.8pt, arc=3pt},
    caja_alerta/.style={colback=orange!5, colframe=orange!80!black, fonttitle=\bfseries, boxrule=0.8pt, arc=3pt}
}

\lstset{
    language=Python,
    basicstyle=\ttfamily\footnotesize,
    keywordstyle=\color{blue}\bfseries,
    commentstyle=\color{codegreen}\itshape,
    stringstyle=\color{codepurple},
    numbers=left, numberstyle=\tiny\color{gray},
    frame=single, backgroundcolor=\color{gray!5},
    breaklines=true, showstringspaces=false
}

% --- Estilo de Página ---
\pagestyle{fancy}
\fancyhf{}
\lhead{\textcolor{ucbblue}{\textbf{IMT-121 | Guía Maestra y Kit de Ejecución}}}
\rhead{\textbf{Sesión 07 -- Semana 3}}
\cfoot{Página \thepage}
\renewcommand{\headrulewidth}{0.4pt}
\renewcommand{\headrule}{\hbox to\headwidth{\color{ucbblue}\leaders\hrule height \headrulewidth\hfill}}

\begin{document}

\begin{center}
    {\LARGE \textbf{GUÍA MAESTRA Y KIT DE EJECUCIÓN}} \\[0.2cm]
    {\Large \textbf{SESIÓN 07}} \\[0.4cm]
    \normalsize
    \textbf{Asignatura:} IMT-121 Circuitos Electrónicos I \\
    \textbf{Carrera:} Ingeniería Mecatrónica e Ingeniería Biomédica \\
    Universidad Católica Boliviana ``San Pablo'' — Sede Tarija \\[0.2cm]
    \textbf{Docente:} M.Sc. Ing. Bernardo Quiroga Turdera \\
    \textbf{Gestión:} Semestre II/2026 \hspace{0.5cm} \textbf{Semana 3 — Sesión 2} (Miércoles 19 de Agosto de 2026) \\[0.2cm]
    \textit{Marco Pedagógico: Criterios ABET (1 y 6) | Metodología CDIO (Concebir–Diseñar–Implementar) | Formulación MNA}
    \vspace{0.3cm}
\end{center}

\hrule
\vspace{0.4cm}

\section*{📌 Diagnóstico y Justificación del Reajuste Curricular}
\begin{tcolorbox}[caja_diagnostico, title=Transición Metodológica]
La secuencia didáctica no está desalineada; se encuentra en el punto exacto de transición metodológica. En las sesiones previas se abordó la formulación matricial por inspección directa para redes pasivas estándar ($\mathbf{R}\mathbf{I}=\mathbf{V}$ y $\mathbf{G}\mathbf{V}=\mathbf{I}$) y el lunes se ejecutó el Laboratorio N$^\circ$ 1 (validación en protoboard con instrumentos reales).

La inspección directa clásica colapsa ante dos elementos indispensables en mecatrónica y biomédica:
\begin{enumerate}
    \item Fuentes de tensión conectadas entre nodos no referenciados (Supernodos).
    \item Fuentes dependientes / controladas (VCVS, VCCS, CCVS, CCCS), que representan los modelos lineales de transductores activos, amplificadores de instrumentación y transistores.
\end{enumerate}
El ajuste para la sesión de hoy consiste en transformar la clase en un \textbf{Taller Teórico-Computacional de Análisis Nodal Modificado (MNA)}, que constituye el núcleo matemático interno de simuladores como SPICE/LTspice, complementándolo con un taller práctico de simulación avanzada (directivas \texttt{.tf} y \texttt{.step}).
\end{tcolorbox}

\section{I. Guion de Clase Minuto a Minuto (90 Minutos)}
\begin{center}
    \renewcommand{\arraystretch}{1.3}
    \begin{tabularx}{\textwidth}{|>{\columncolor{gray!10}\bfseries}c|X|}
        \hline
        \multicolumn{2}{|c|}{\cellcolor{ucbblue}\textcolor{white}{\textbf{ESTRUCTURA DE LA SESIÓN 07 (90 MIN)}}} \\
        \hline
        00--10 min & Debrief del Laboratorio 1 (Discusión de $\varepsilon_r$ y efecto carga) \\
        \hline
        10--35 min & Bloque Teórico: Supernodos y Fuentes Dependientes en MNA \\
        \hline
        35--55 min & Ejemplo Integrador en Pizarra (Cálculo Analítico Completo) \\
        \hline
        55--75 min & Taller de Simulación en LTspice: Directivas \texttt{.tf} y \texttt{.step} \\
        \hline
        75--85 min & Live-Coding en Python (Resolución de matrices aumentadas) \\
        \hline
        85--90 min & Ejercicio de Cierre en Tríadas y Asignación de Estudio \\
        \hline
    \end{tabularx}
\end{center}

\section{II. Fundamentación Teórica: El Método Nodal Modificado (MNA)}
\subsection{1. Limitación de la Inspección Directa}
Cuando una red contiene fuentes de voltaje (independientes o dependientes) conectadas entre dos nodos esenciales \textit{sin} resistencia en serie, no es posible definir directamente la corriente que fluye por ellas mediante la Ley de Ohm ($\Delta V / R$ con $R \to 0$).

\subsection{2. Tratamiento de Supernodos}
Un Supernodo se forma al englobar una fuente de voltaje conectada entre dos nodos no referenciados y cualquier elemento conectado en paralelo con ella.

\begin{center}
    \begin{circuitikz}[scale=1, transform shape]
        \draw (0,0) node[left]{Nodo 2} to[short, o-] (1.5,0)
              to[V, v=$V_d$] (3.5,0)
              to[short, -o] (5,0) node[right]{Nodo 3};
        % Caja del supernodo
        \draw[dashed, red, thick, rounded corners] (0.5, -1) rectangle (4.5, 1);
        \node[red, above] at (2.5, 1) {\textbf{SUPERNODO (2-3)}};
    \end{circuitikz}
\end{center}

\textbf{Ecuación de Conservación (LCK):} Se aplica la Ley de Corrientes de Kirchhoff a la superficie cerrada del supernodo completo:
\begin{equation}
    \sum I_{\text{salientes de (Nodo 2)}} + \sum I_{\text{salientes de (Nodo 3)}} = 0
\end{equation}

\textbf{Ecuación de Enlace / Restricción:} Se define la diferencia de potencial entre los terminales de la fuente:
\begin{equation}
    V_2 - V_3 = V_{\text{fuente}}
\end{equation}

\subsection{3. Incorporación de Fuentes Dependientes en el Sistema Matricial}
Las fuentes controladas introducen términos que rompen la simetría de la matriz de conductancias $\mathbf{G}$. Para resolver el sistema:
\begin{itemize}
    \item Se expresan las variables de control (ej. voltajes diferenciales $V_x$ o corrientes de rama $I_x$) en función exclusiva de los potenciales nodales incógnita ($V_1, V_2, \dots, V_N$).
    \item Se reagrupan los coeficientes en el miembro izquierdo de la ecuación matricial $\mathbf{A}\mathbf{x} = \mathbf{b}$.
\end{itemize}

\section{III. Ejemplo Integrador Resuelto Paso a Paso (Guía de Estudio)}
\textbf{Problema de Ingeniería:} \\
El siguiente circuito modela la etapa de preamplificación de un sensor acoplado. Contiene una fuente de polarización $V_s = \qty{10}{\volt}$, resistencias de acoplamiento y una fuente de tensión controlada por voltaje (VCVS) que actúa como buffer con ganancia dependiente: $V_d = 3 V_x$, donde el voltaje de control es $V_x = V_2$.

\begin{center}
    \begin{circuitikz}[scale=1, transform shape]
        % Nodos y ramas principales
        \draw (0,0) to[V, v=$V_s {=} \qty{10}{\volt}$] (0,3) node[above]{Nodo 1}
              to[R, l=$R_1 {=} \qty{50}{\ohm}$] (3,3) node[above]{Nodo 2}
              to[cV, v=$3V_x$] (6,3) node[above]{Nodo 3}
              to[R, l=$R_3 {=} \qty{200}{\ohm}$] (6,0) -- (0,0);
        \draw (3,3) to[R, l=$R_2 {=} \qty{100}{\ohm}$] (3,0);
        % Tierras
        \draw (0,0) node[ground]{};
        \draw (3,0) node[ground]{};
        \draw (6,0) node[ground]{};
    \end{circuitikz}
\end{center}
\textit{(Nota de topología: $V_s$ fija el Nodo 1 a $\qty{10}{\volt}$. Entre el Nodo 2 y el Nodo 3 se encuentra la fuente dependiente $3V_x$, formando un Supernodo).}

\vspace{0.3cm}
\textbf{Paso 1: Identificación de Nodos y Ecuación de Control} \\
Nodo 1 está conectado directamente a la fuente independiente:
\begin{equation}
    V_1 = \qty{10.0}{\volt}
\end{equation}
Variable de Control:
\begin{equation}
    V_x = V_2
\end{equation}

\textbf{Paso 2: Ecuación de Enlace del Supernodo (2-3)} \\
La fuente dependiente impone la diferencia de potencial entre el Nodo 2 y el Nodo 3:
\begin{equation}
    V_2 - V_3 = 3 V_x = 3 V_2
\end{equation}
Reagrupando términos obtenemos la restricción estructural:
\begin{equation}
    V_2 - 3 V_2 - V_3 = 0 \implies \mathbf{-2 V_2 - V_3 = 0} \implies V_3 = -2 V_2
\end{equation}

\textbf{Paso 3: LCK en el Supernodo (Caso Flotante y Ajuste de Topología Realista)}
\begin{tcolorbox}[caja_alerta, title=Discusión de Ingeniería en Pizarra]
Si asumiéramos topológicamente que \textit{no} existe $R_1$ (es decir, $R_1 \to 0$ y $V_s$ conecta directo a $R_2$ vía una red $\qty{100}{\ohm}$ como en un boceto inicial simplificado), la ecuación nodal sería $\frac{V_2 - 10}{100} + \frac{V_3}{200} = 0 \implies 2V_2 + V_3 = 20$. 
Al cruzarla con $-2V_2 - V_3 = 0$, obtenemos $0 = 20$, una \textbf{inconsistencia física}. \\
\textbf{¿Por qué ocurre?} Porque la VCVS intenta imponer una ganancia infinita sin rama de atenuación/realimentación de entrada. Al incluir $R_1 = \qty{50}{\ohm}$, el circuito se vuelve estable.
\end{tcolorbox}

Aplicamos LCK al contorno cerrado (2-3) con la red completa:
\begin{itemize}
    \item Corriente saliendo del Nodo 2 hacia el Nodo 1: $\frac{V_2 - V_1}{R_1} = \frac{V_2 - 10}{50}$
    \item Corriente saliendo del Nodo 2 hacia GND: $\frac{V_2}{R_2} = \frac{V_2}{100}$
    \item Corriente saliendo del Nodo 3 hacia GND: $\frac{V_3}{R_3} = \frac{V_3}{200}$
\end{itemize}
Ecuación nodal del Supernodo:
\begin{equation}
    \frac{V_2 - 10}{50} + \frac{V_2}{100} + \frac{V_3}{200} = 0
\end{equation}
Multiplicando por $200$ para eliminar denominadores:
\begin{equation}
    4(V_2 - 10) + 2 V_2 + V_3 = 0 \implies \mathbf{6 V_2 + V_3 = 40}
\end{equation}

\textbf{Paso 4: Resolución del Sistema $2\times 2$}
Combinando la ecuación del supernodo con la restricción ($V_3 = -2 V_2$):
\begin{align*}
    6 V_2 + (-2 V_2) &= 40 \\
    4 V_2 &= 40 \implies \mathbf{V_2 = \qty{10.0}{\volt}}
\end{align*}
Evaluando los demás potenciales:
\begin{align*}
    \mathbf{V_3} &= -2(10.0) = \mathbf{\qty{-20.0}{\volt}} \\
    \mathbf{V_x} &= \mathbf{\qty{10.0}{\volt}}
\end{align*}

\section{IV. Cómputo Científico en Python (NumPy)}
Script para resolver sistemas nodales con fuentes controladas mediante matrices aumentadas:

\begin{lstlisting}[language=Python]
import numpy as np

# 1. Definicion del sistema A * V = b
# Variables: [V2, V3]
# Ecuacion 1 (Supernodo LCK):  6*V2 + 1*V3 = 40
# Ecuacion 2 (Restriccion VCVS): -2*V2 - 1*V3 = 0

A = np.array([
    [ 6.0,  1.0],
    [-2.0, -1.0]
], dtype=float)

b = np.array([40.0, 0.0], dtype=float)

# 2. Calculo matricial
det_A = np.linalg.det(A)
assert np.abs(det_A) > 1e-6, "Matriz singular: Revise dependencias de fuentes."

V_sol = np.linalg.solve(A, b)
V2, V3 = V_sol[0], V_sol[1]

print("--- RESULTADOS MATRICIALES (PYTHON) ---")
print(f"Determinante del sistema: {det_A:.2f}")
print(f"Voltaje V2 (Control Vx) : {V2:.4f} V")
print(f"Voltaje V3 (Salida)     : {V3:.4f} V")
\end{lstlisting}

\newpage
\section{V. Taller de Simulación Avanzada en LTspice (Análisis \texttt{.tf})}
En lugar de limitarse al análisis de punto de operación (\texttt{.op}), los estudiantes aprenderán a utilizar la directiva de Función de Transferencia en CC (\texttt{.tf}), empleada en la industria para calcular automáticamente:
\begin{itemize}[label=\textcolor{ucbblue}{\textbullet}, itemsep=0pt]
    \item La ganancia de pequeña señal ($\frac{V_{\text{out}}}{V_{\text{in}}}$).
    \item La impedancia de entrada de la red ($R_{\text{in}}$).
    \item La resistencia de salida o Resistencia de Thévenin ($R_{\text{out}} = R_{th}$).
\end{itemize}

\begin{center}
    \begin{circuitikz}[scale=0.9, transform shape]
        \draw (0,0) to[V, v=\texttt{V1 10}] (0,3)
              to[R, l=\texttt{R1 50}] (3,3) node[above]{\texttt{N002}}
              to[cV, v=\texttt{E1 Gain=3}] (6,3) node[above]{\texttt{N003}}
              to[R, l=\texttt{R3 200}] (6,0) -- (0,0);
        \draw (3,3) to[R, l=\texttt{R2 100}] (3,0);
        \draw (0,0) node[ground]{}; \draw (3,0) node[ground]{}; \draw (6,0) node[ground]{};
    \end{circuitikz}
\end{center}

\textbf{Instrucciones para el Estudiante en Laboratorio de Software:}
\begin{enumerate}
    \item \textbf{Colocar la Fuente Controlada:} Presionar \tecla{F2} y seleccionar el componente \texttt{e} (Voltage Dependent Voltage Source - VCVS).
    \item \textbf{Conexión de Pines de E:}
    \begin{itemize}[label=$\circ$]
        \item Pines de salida: conectarlos entre el Nodo 2 y el Nodo 3.
        \item Pines de control (+,-): conectarlos al Nodo 2 y a Tierra (para medir $V_x = V_2$).
    \end{itemize}
    \item \textbf{Asignar el valor de ganancia:} Hacer clic derecho y asignar \texttt{Value = 3}.
    \item \textbf{Directiva SPICE:} Presionar \tecla{.} e ingresar: \texttt{.tf V(N003) V1}
    \item \textbf{Ejecutar Simulación:} LTspice entregará una ventana con:
    \begin{itemize}
        \item \texttt{Transfer\_function = -2.000} \quad $\left(\frac{V_3}{V_1} = \frac{\qty{-20}{\volt}}{\qty{10}{\volt}} = -2\right)$
        \item \texttt{v1\#Input\_impedance = 83.333 ohms}
        \item \texttt{output\_impedance\_at\_v(n003) = 0 ohms} (característico de una fuente de tensión ideal).
    \end{itemize}
\end{enumerate}

\section{VI. Ejercicio Activo en Tríadas (Trabajo en Aula — 15 MIN)}
\textbf{Enunciado:} \\
Para el circuito analizado, modifique la ganancia de la fuente controlada $E_1$ a un valor de $k = 1.5$ ($V_d = 1.5 V_x$).
\begin{enumerate}
    \item Deduzca analíticamente el nuevo voltaje $V_2$ y $V_3$.
    \item Modifique el parámetro en su script de Python y compruebe el resultado.
    \item Modifique el valor en LTspice y confirme la nueva función de transferencia $\frac{V_3}{V1}$.
\end{enumerate}

\begin{tcolorbox}[caja_diagnostico, title=Solución de Autocontrol Docente]
Con la nueva restricción $V_3 = (1 - 1.5)V_2 = -0.5 V_2$, al sustituir en la ecuación del supernodo: \\
$6 V_2 + V_3 = 40 \implies 6 V_2 - 0.5 V_2 = 40 \implies 5.5 V_2 = 40$ \\
$V_2 = \mathbf{\qty{7.273}{\volt}}$ \quad|\quad $V_3 = \mathbf{\qty{-3.636}{\volt}}$ \quad|\quad Ganancia = \textbf{-0.3636}
\end{tcolorbox}

\vspace{0.3cm}
\textit{Al finalizar esta sesión, los estudiantes dominarán la formulación analítica de circuitos no triviales con fuentes activas y sabrán cómo validar modelos lineales complejos empleando Python y LTspice de manera concurrente.}

\end{document}